
\subsubsection{6.1 Interpretation of Results}

\textbf{Finding 1: Cluster-group alignment depends on pretrained initialization alignment with the spurious attribute.}

Both Waterbirds and CelebA produce spurious correlation conditions under which ERM models exhibit degraded worst-group accuracy, yet only Waterbirds yields cluster-separable spurious groups at epoch 5. The differentiating factor is the representation encoded by the ImageNet-pretrained ResNet-50 prior to any task-specific fine-tuning. ImageNet training produces representations that encode bird habitats, scene-level backgrounds, and environmental contexts as high-salience, separable semantic dimensions. These features align directly with the Waterbirds spurious attribute. By contrast, hair color is not a primary semantic category in ImageNet; the pretrained model does not pre-separate samples along a hair-color axis. Five epochs of ERM fine-tuning on CelebA are insufficient to construct hair-color-separable cluster structure from this starting point.

An alternative partial explanation is that k=2 underspecifies the CelebA group structure. CelebA has four groups (binary class × binary spurious attribute), with the blonde-male group comprising approximately 1\% of training data. k=2 clustering collapses this structure and cannot simultaneously recover both class and spurious axes. This underspecification contributes to the CelebA failure alongside the pretraining alignment explanation; separating the contributions of these two factors would require a k=4 ablation, which was not conducted.

\textbf{Finding 2: The CelebA failure is near-random, not marginal.}

The difference between the CelebA purity (0.456) and the random baseline (0.440) is 0.016. This margin is below the resolution that would be considered informative for practical purposes. A practitioner deploying epoch-5, k=2 clustering on CelebA embeddings would obtain cluster assignments that carry essentially no spurious group signal above chance. Any downstream robustification using these assignments would be statistically equivalent to random reweighting, without any indication to the practitioner that detection has failed.

\textbf{Finding 3: Published benchmark results do not transfer across experimental configurations.}

The configuration sensitivity documented here — early-epoch versus convergence-epoch, k=2 versus k≥4 — is practically significant for intervention methods that require detection at early training epochs. Reporting clustering results obtained under post-convergence, flexible-k conditions as evidence of detection reliability in early-epoch, k=2 settings is not well-supported by these data.

\subsubsection{6.2 Proposed Diagnostic: Pretrained Linear Probe Pre-Screen}

Based on the interpretation that detection success depends on pretraining alignment, the following diagnostic is proposed as a candidate gate before deploying epoch-5 clustering-based spurious detection:

Extract penultimate-layer embeddings from the pretrained model prior to any task-specific fine-tuning. Train a linear classifier on a small labeled set for the suspected spurious attribute using these pretrained embeddings. If linear probe accuracy exceeds a heuristic threshold (approximately 70\% is tentatively suggested based on the interpretation that this indicates the spurious attribute is encoded as a linearly separable feature in the pretrained representation), epoch-5 k-means detection is likely to produce above-chance cluster-group alignment. If accuracy is near chance, early-epoch k-means will likely fail.

This pre-screen requires spurious-attribute labels, not full group labels. Spurious-attribute labels are generally less expensive to collect than group labels. The 70\% threshold is a heuristic that requires empirical calibration across datasets and architectures; the present work does not provide this calibration. Validation of the diagnostic, including threshold determination and cross-architecture testing, is necessary before it can be applied reliably.

\subsubsection{6.3 Limitations}

\textbf{L1 — Single random seed.} All experiments use one fixed seed (seed=42 for clustering, seed=1 for training). The k-means initialization variance is partially mitigated by n_init=10. ERM training seed variance is not characterized. The Waterbirds AMI margin above threshold is 0.262 (AMI=0.762 versus threshold=0.500), which is large relative to typical k-means variance, but formal characterization across ≥5 seeds is absent.

\textbf{L2 — k=2 underspecification for multi-group structure.} CelebA has a 4-group spurious structure (binary class × binary spurious attribute). k=2 conflates this structure, particularly affecting the minority blonde-male group (approximately 1\% of training samples). Disambiguating the contribution of k=2 underspecification from pretraining alignment as explanations for the CelebA failure requires a k=4 ablation. This ablation was not performed.

\textbf{L3 — Probe epoch not normalized across datasets.} Waterbirds has 4,795 training samples; CelebA has 162,770. Epoch 5 represents substantially different amounts of training in terms of total gradient updates. A principled probe epoch selection criterion based on representational stability (e.g., cosine similarity between consecutive epoch embeddings) was not implemented.

\textbf{L4 — Downstream mechanism not evaluated.} The Gradient SNR Balancing (GSB) intervention — equalizing gradient SNR along cluster-discovered directions to improve worst-group accuracy — was not evaluated. The prerequisite detection condition (H-E1) failed on CelebA, blocking all downstream hypothesis experiments (H-M1 through H-C1). The entire mechanistic chain from gradient SNR imbalance to representational suppression of invariant features remains theoretically motivated but empirically unvalidated.

\textbf{L5 — PruSC citation not independently verified.} The attribution of >95\% CelebA purity to Kim et al. [2024] (PruSC) was not independently confirmed through a literature search. This citation was sourced from internal pipeline documentation (045_validated_hypothesis.md). If the cited paper does not exist or does not report this value, the comparison in Section 5.5 is unsupported.

\textbf{L6 — Single architecture and modality.} Experiments are conducted with ResNet-50 on two visual classification datasets. Generalization to other architectures (e.g., ViT-B) and other modalities (text, audio) is not established.

\subsubsection{6.4 Broader Considerations}

Raising awareness of the feature-strength conditionality of annotation-free detection is intended to encourage validation of the detection step before trusting downstream robustification results. The proposed pre-screen provides a candidate method for this validation, pending empirical confirmation. There are no identified direct negative societal impacts from this diagnostic characterization work.

